\begin{example}[SPARQL query in a tensor network]
    Let us consider the example query:
    \begin{minted}{sparql}
    PREFIX foaf: <http://xmlns.com>
    SELECT ?a ?c
    WHERE{
        ?a foaf:knows ?b;
        ?b foaf:friendOf ?c;
    }
    \end{minted}

    As a tensor network the query is represented as:
    \stantikzpic{

        \drawvwire{-2.5}{-1}{$\indvariableof{a}$}
        \drawtensorblock{0}{0}{3}{$\kg$}
        \drawvwire{0}{-1}{}
        \drawtensorblock{0}{-4}{1.5}{$\onehotmapof{knows}$}

        \draw (3.5,-2.5) to[bend right= -10] (2.5,-1);
        \drawvariabledot{3.5}{-2.5}
        \drawtextnode{3.5}{-3.25}{\colorlabelsize $\indvariableof{b}$}
        \draw (3.5,-2.5) to[bend right= 10] (4.5,-1);

        \drawtensorblock{7}{0}{3}{$\kg$}
        \drawvwire{7}{-1}{}
        \drawtensorblock{7}{-4}{2}{$\onehotmapof{friendOf}$}
        \drawvwire{9.5}{-1}{$\indvariableof{c}$}
    }


    We can now apply sparsification techniques developed in \secref{sec:groundingStorage} to provide an efficient storage.
    \begin{itemize}
        \item Local contraction with the basis vectors:
        \stantikzpic{

            \drawtensorblock{-9}{0}{2}{$\hypercoreof{knows}$}
            \drawvwire{-10.5}{-1}{$\indvariableof{a}$}
            \drawvwire{-7.5}{-1}{$\indvariableof{b}$}

            \drawtextnode{-5}{-1}{$\algdefsymbol$}

            \drawvwire{-2.5}{-1}{$\indvariableof{a}$}
            \drawtensorblock{0}{0}{3}{$\kg$}
            \drawvwire{0}{-1}{}
            \drawtensorblock{0}{-4}{1.5}{$\onehotmapof{knows}$}
            \drawvwire{-2.5}{-1}{$\indvariableof{a}$}
            \drawvwire{2.5}{-1}{$\indvariableof{b}$}

            \drawtextnode{5}{-1}{and}

            \begin{scope}[shift={(20,0)}]

                \drawtensorblock{-9}{0}{2.5}{$\hypercoreof{friendsOf}$}
                \drawvwire{-10.5}{-1}{$\indvariableof{b}$}
                \drawvwire{-7.5}{-1}{$\indvariableof{c}$}

                \drawtextnode{-5}{-1}{$\algdefsymbol$}

                \drawtensorblock{0}{0}{3}{$\kg$}
                \drawvwire{0}{-1}{}
                \drawtensorblock{0}{-4}{2}{$\onehotmapof{firendsOf}$}
                \drawvwire{-2.5}{-1}{$\indvariableof{b}$}
                \drawvwire{2.5}{-1}{$\indvariableof{c}$}
            \end{scope}
        }
        \item Message passing of the support:
        \begin{align*}
            \sechypercoreofat{knows}{\indvariableof{a},\indvariableof{b}}
            \algdefsymbol \nzcontractionof{
                \hypercoreofat{knows}{\indvariableof{a},\indvariableof{b}},
                \contractionof{\hypercoreofat{friendsOf}{\indvariableof{b},\indvariableof{c}}}{\indvariableof{b}}
            }{\indvariableof{a},\indvariableof{b}}
        \end{align*}
        and
        \begin{align*}
            \sechypercoreofat{friendsOf}{\indvariableof{b},\indvariableof{c}}
            \algdefsymbol \nzcontractionof{
                \hypercoreofat{friendsOf}{\indvariableof{b},\indvariableof{c}},
                \contractionof{\hypercoreofat{knows}{\indvariableof{a},\indvariableof{b}}}{\indvariableof{b}}
            }{\indvariableof{b},\indvariableof{c}}
        \end{align*}
    \end{itemize}
    The resulting sparsified storage of the query is then the tensor network
    \begin{align*}
        \{\sechypercoreofat{knows}{\indvariableof{a},\indvariableof{b}},\sechypercoreofat{friendsOf}{\indvariableof{b},\indvariableof{c}}\} \, .
    \end{align*}
    For any locally supported index, we find an extension to an global index which is supported by the tensor network (see \charef{cha:messagePassing}).
    For example, if for a tuple $\indindexof{a},\indindexof{b}$ we have $\sechypercoreofat{knows}{\indexedindvariableof{a},\indexedindvariableof{b}}=1$ then there is $\indindexof{c}$ with $\sechypercoreofat{friendsOf}{\indexedindvariableof{b},\indexedindvariableof{c}}=1$, which means that $(\indindexof{a},\indindexof{b},\indindexof{c})$ is a solution of the query.
\end{example}